Visual world models for control have progressed from pixel reconstruction
to joint-embedding prediction, enabling latent-space planning with
Cross-Entropy Method (CEM) and related MPC-style search.
Yet most JEPA-style stacks still advance latent state with a
\emph{discrete} autoregressive (AR) next-embedding map: one nonlinear
step must absorb an entire physical transition, with no explicit flow
along the planning interval and limited inductive bias toward temporal
coherence.
We propose \method{}, a visual world model whose latent dynamics are
\emph{continuous}: an action-conditioned neural ODE
$\mathrm{d}z/\mathrm{d}\tau = v_\theta(z,a)$, with one planning step
defined as fixed-step RK4 integration from $\tau{=}0$ to $\tau{=}1$.
Training keeps the standard JEPA next-embedding objective together with
SIGReg isotropic-Gaussian regularization; encoder, projectors, and the
CEM loop are unchanged so comparisons isolate continuous vs.\ discrete
latent evolution.
On the official PushT expert dataset with a fixed 1000-window subset
(seed~3072; train~901 / val~99), \method{} matches the discrete AR
baseline on val-tied CEM success (\textbf{5/93, $\approx$5.38\%}) while
achieving lower train and validation losses (0.234 / 0.506 vs.\
0.248 / 0.536).
We report limitations honestly: this subset is far below full-data
regimes that report high PushT success, and CEM is tied despite the
loss gain.
All numbers use the official Hugging Face release; no synthetic episodes
are used.
